\section{Проектирование приложения}

\subsection{Варианты использования системы}

Перед описанием архитектурных решений целесообразно зафиксировать состав пользователей системы и крупные блоки доступных им функций.
В системе выделены пять ролей, образующих иерархию по объёму прав на запись: Просматривающий (\code{DataViewer}), Оператор приёмной комиссии (\code{AdmissionsOperator}), Аудитор (\code{Auditor}), Администратор данных (\code{DataAdministrator}) и Суперадминистратор (\code{SuperAdmin}).
Все аутентифицированные пользователи имеют право просмотра всех данных; различия касаются только прав на запись.
Иерархия отражает накопление прав: каждая следующая роль обладает всеми возможностями предыдущей и дополнительной группой функций.

Диаграмма вариантов использования представлена на рисунке~\ref{fig:usecases}.
На ней показаны крупные категории сценариев, объединяющие однородные операции (детализация конкретных операций каждой категории даётся в подразделах ниже и в разделе~5).

\begin{figure}[htbp]
\centering
\includegraphics[width=\textwidth]{use-cases}
\caption{Диаграмма вариантов использования}
\label{fig:usecases}
\end{figure}

Соответствие категорий сценариев и ролей следующее.
Сценарий <<Аутентификация>> доступен всем пользователям, включая неаутентифицированных, и предшествует любой работе с системой.
Сценарий <<Просмотр результатов отбора и выгрузка в Excel>> и сценарий <<Просмотр журналов действий и входов>> доступны всем аутентифицированным пользователям, начиная с роли Просматривающего.
Сценарий <<Ведение данных абитуриентов и заявок>> доступен Оператору приёмной комиссии и всем ролям выше по иерархии.
Сценарий <<Валидация данных и разрешение удалений>> доступен Аудитору и всем ролям выше.
Сценарий <<Настройка правил приёма и запуск пересчёта>> объединяет управление организационной структурой, конкурсными списками, категориями приёма и группами критериев оценивания; он доступен Администратору данных и Суперадминистратору.
Сценарий <<Управление учётными записями пользователей>> доступен только Суперадминистратору.

Подробное соответствие ролей и прав на запись приведено в таблице~\ref{tab:roles} подраздела~\ref{subsec:roles}.

\subsection{Общая структура взаимодействия компонентов}

Для системы выбрана клиент-серверная архитектура с разделением обязанностей между клиентским одностраничным приложением (Single-Page Application), серверной частью REST API и реляционной базой данных.

Клиентская часть~--~одностраничное приложение, выполняемое в браузере пользователя.
Клиент взаимодействует с сервером исключительно через REST API по протоколу HTTPS, передавая данные в формате JSON.
Для аутентификации пользователя выбрана схема Bearer-токена: после успешного входа клиент сохраняет полученный JSON Web Token в локальном хранилище браузера и автоматически добавляет его в заголовок каждого последующего запроса.

Серверная часть выполняет роль REST API: принимает HTTP-запросы от клиента, выполняет проверку прав доступа, валидацию входных данных, реализует бизнес-логику и взаимодействует с базой данных.
Серверная часть не хранит состояния сессии пользователя~--~вся необходимая информация о пользователе извлекается из проверенного JWT-токена при каждом запросе.
Подобный подход обеспечивает горизонтальную масштабируемость и упрощает развёртывание системы.

База данных обеспечивает долговременное хранение всех данных приложения с поддержкой транзакционной целостности на уровне СУБД.
Управление схемой базы данных возлагается на отдельную систему миграций, что позволяет отслеживать историю изменений структуры и автоматически применять необходимые миграции при обновлении приложения.

Схема общего взаимодействия компонентов приложения представлена на рисунке~\ref{fig:architecture}.

\begin{figure}[htbp]
\centering
\includegraphics[width=0.9\textwidth]{architecture}
\caption{Схема взаимодействия компонентов приложения}
\label{fig:architecture}
\end{figure}

\subsection{Архитектура серверной части}

Архитектура серверной части~--~двухслойная и построена по шаблону <<репозиторий>> (Repository) из каталога шаблонов корпоративных приложений М.~Фаулера \citesrc{fowler-paa}: каждой таблице базы данных соответствует отдельный класс-репозиторий, инкапсулирующий доступ к ней, а контроллеры обращаются к репозиториям напрямую, без выделения отдельного сервисного слоя.
Подобный подход выбран сознательно: большинство операций над данными представляют собой стандартные операции CRUD, для которых введение дополнительного абстрактного слоя избыточно и приводит к разрастанию количества классов без увеличения функциональной ценности.
Бизнес-логика, выходящая за рамки CRUD (алгоритм отбора, экспорт в Excel, ведение журнала аудита), выносится в отдельные сервисы, подключаемые к репозиториям и контроллерам.

Слой контроллеров принимает HTTP-запросы, выполняет преобразование входных данных в объекты передачи данных, декларативно проверяет права доступа по роли пользователя, делегирует выполнение операции репозиторию и возвращает HTTP-ответ.
Слой репозиториев инкапсулирует доступ к базе данных, выполняет преобразование данных между типами языка реализации и типами СУБД, реализует логику CRUD-операций и вызывает сервисы аудита и пересчёта результатов отбора.

\subsection{Архитектура клиентской части}

Клиентская часть строится на основе библиотеки React с применением функциональных компонентов и хуков.
Маршрутизация выполняется средствами библиотеки декларативной маршрутизации, что позволяет компактно описать карту страниц и ограничивать доступ к ним в зависимости от роли текущего пользователя.

Глобальное состояние клиентского приложения разделяется на три независимые области: сведения о текущем пользователе и его токене доступа, текущая тема оформления (светлая или тёмная) и текущий язык интерфейса.
Каждая из областей инкапсулируется в отдельный контекст и не зависит от остальных.
Состояние, специфичное для конкретной страницы (фильтры таблицы, открытые модальные окна), хранится в локальном состоянии соответствующего компонента.

Взаимодействие с сервером выполняется через единый клиент HTTP-запросов с настроенным базовым адресом API.
На клиент устанавливаются перехватчики запросов и ответов, обеспечивающие: автоматическое добавление заголовка \code{Authorization: Bearer <token>} ко всем исходящим запросам; добавление заголовка \code{Accept-Language} с текущим языком интерфейса; обработку ошибки 401 (Unauthorized) с автоматическим переходом на страницу входа.

\subsection{Проектирование базы данных}

\subsubsection{Метод проектирования}

Для проектирования базы данных выбран классический подход <<сверху вниз>>: на основе анализа предметной области строится концептуальная модель данных, которая последовательно детализируется до логической и физической моделей.
Схема базы данных описывается набором SQL-сценариев, объединённых в один changelog-файл формата YAML под управлением системы миграций.
Подобный подход обеспечивает полный контроль над структурой и индексами базы данных, делает все изменения схемы явными и отслеживаемыми в системе контроля версий, а также гарантирует идентичность схемы в средах разработки, тестирования и эксплуатации.

\subsubsection{Концептуальная модель}

В концептуальной модели выделены следующие группы сущностей.

Организационная иерархия: факультет содержит кафедры, кафедра содержит специальности.
Для каждой специальности определены один или несколько конкурсных списков (бюджет, платное, целевое направление), каждый конкурсный список разделён на категории приёма с указанием квоты и приоритета защиты.

Критерии оценивания: для каждого конкурсного списка определена группа критериев оценивания (например, средний балл аттестата, баллы централизованного тестирования по предметам).
Группа состоит из критериев, упорядоченных по приоритету для лексикографической сортировки кандидатов.

Абитуриенты: каждый абитуриент имеет идентификатор, имя, оценки по критериям и список заявок в категории приёма с указанием приоритета предпочтения.
Для каждого абитуриента сохраняется булевый признак валидации данных вторым пользователем.

Результаты отбора: итоговая таблица содержит идентификаторы абитуриентов и категорий, в которые они зачислены.
Таблица полностью перестраивается алгоритмом пересчёта при любом изменении входных данных.

Подсистемы аудита и логического удаления: журнал действий пользователей, журнал попыток входа в систему, запросы на удаление абитуриентов с двухступенчатым подтверждением.

ER-диаграмма базы данных представлена на рисунке~\ref{fig:erdiagram}.

\begin{figure}[htbp]
\centering
\includegraphics[width=\linewidth,keepaspectratio]{er}
\caption{ER-диаграмма базы данных}
\label{fig:erdiagram}
\end{figure}

\subsubsection{Логическая модель}

Логическая модель базы данных нормализована до третьей нормальной формы.
Ключевые таблицы и их связи приведены в таблице~\ref{tab:tables}.

\begin{longtblr}[
  caption = {Основные таблицы базы данных},
  label = {tab:tables},
]{
  colspec = {|Q[l]|X|},
  width = \textwidth,
  hlines, vlines,
  rowhead = 1,
  row{1} = {halign=c},
}
Таблица & Назначение \\
\code{faculties} & Факультеты университета \\
\code{departments} & Кафедры, связь с факультетом по внешнему ключу \\
\code{specialties} & Специальности, связь с кафедрой по внешнему ключу \\
\code{competitionlists} & Конкурсные списки специальности с общим планом приёма \\
\code{admissioncategories} & Категории приёма с квотой и приоритетом защиты \\
\code{evaluationcriteria} & Критерии оценивания с диапазоном и выбором варианта <<больше~--~лучше>> или <<меньше~--~лучше>> \\
\code{evaluationcriteriagroups} & Группы критериев \\
\code{evaluationcriteriagroupitems} & Связь группы и критерия с приоритетом для сортировки \\
\code{applicants} & Абитуриенты с признаком валидации данных \\
\code{applicantevaluationvalues} & Оценки абитуриентов по критериям \\
\code{applicantadmissioncategories} & Заявки абитуриентов в категории с приоритетом предпочтения \\
\code{selectedapplicants} & Итоговые результаты отбора \\
\code{users} & Учётные записи операторов с ролью и признаком активности \\
\code{audit\_log} & Журнал действий пользователей \\
\code{auth\_log} & Журнал попыток входа в систему \\
\code{applicant\_deletion\_requests} & Запросы на удаление абитуриентов \\
\end{longtblr}

Для таблицы \code{applicantevaluationvalues} предусмотрено уникальное ограничение по паре столбцов \code{(applicant\_id, evaluation\_criteria\_id)}, что гарантирует невозможность хранения двух оценок одного абитуриента по одному критерию.
В таблице \code{applicants} предусмотрен булевый столбец \code{validated} для хранения статуса валидации данных вторым пользователем.

\subsection{Ролевая модель доступа}
\label{subsec:roles}

В системе определены пять ролей пользователей, отличающихся набором прав на запись данных.
Все аутентифицированные пользователи имеют право чтения всех данных.
Соответствие ролей и прав на запись приведено в таблице~\ref{tab:roles}.

\begin{longtblr}[
  caption = {Ролевая модель доступа},
  label = {tab:roles},
]{
  colspec = {|Q[l,wd=4.5cm]|X|},
  width = \textwidth,
  hlines, vlines,
  rowhead = 1,
  row{1} = {halign=c},
}
Роль & Права на запись \\
\code{SuperAdmin} & Все операции, включая управление учётными записями пользователей \\
\code{DataAdministrator} & Структура университета, абитуриенты, аудит, кроме управления пользователями \\
\code{Auditor} & Абитуриенты, валидация данных, подтверждение и отклонение запросов на удаление \\
\code{AdmissionsOperator} & Только данные абитуриентов и их заявок в категории \\
\code{DataViewer} & Только чтение, без прав записи \\
\end{longtblr}

Проверка прав доступа выполняется на серверной части декларативно: каждый метод обработчика помечается требуемой ролью, и попытки вызова метода пользователем с недостаточными правами отклоняются ещё до выполнения бизнес-логики.
На клиентской части все аутентифицированные пользователи видят одинаковый набор страниц, но кнопки выполнения операций скрываются в зависимости от роли текущего пользователя.

\subsection{Алгоритм пересчёта результатов отбора}

Центральным элементом проекта является алгоритм автоматического распределения абитуриентов по категориям приёма.
Алгоритм реализует принцип deferred acceptance \citesrc{gale} с вытеснением слабых кандидатов \citesrc{roth} и состоит из двух фаз: проверка квоты категории и проверка общего плана конкурсного списка.

\subsubsection{Этапы алгоритма}

Этап~1. Предварительная фильтрация.
Из множества заявок исключаются те, для которых у абитуриента не заполнен хотя бы один критерий из группы критериев, привязанной к категории.
Подобные заявки никогда не будут учтены при отборе, поэтому их фильтрация на ранней стадии сокращает объём обрабатываемых данных.

Этап~2. Инициализация.
Все абитуриенты с непустыми списками заявок помещаются в очередь обработки.
Для каждого абитуриента строится упорядоченный по полю \code{selectionpriority} список предпочтений между категориями.

Этап~3. Основной цикл.
Из очереди извлекается очередной абитуриент.
Из его списка предпочтений берётся следующая по приоритету категория.
Абитуриент добавляется в пул кандидатов этой категории.
Пул~--~упорядоченное множество с компаратором, упорядочивающим кандидатов по убыванию <<силы>>.

Этап~4. Проверка квоты категории.
Если после добавления нового кандидата размер пула категории превысил её квоту, из пула извлекается самый слабый кандидат~--~максимальный элемент множества при компараторе по убыванию.
Извлечённый кандидат возвращается в очередь обработки, и его указатель текущего предпочтения смещается на следующую категорию.

Этап~5. Проверка общего плана конкурсного списка.
Если суммарное количество кандидатов во всех категориях одного конкурсного списка превысило его план, из категории с наибольшим значением \code{priority} (наименее защищённой) извлекается самый слабый кандидат и возвращается в очередь.

Этап~6. Завершение.
Алгоритм завершается, когда очередь обработки становится пустой.
Содержимое всех пулов категорий копируется в таблицу \code{selectedapplicants} массовой операцией после полной очистки таблицы.

\subsubsection{Асимптотическая сложность алгоритма}

Каждая заявка обрабатывается не более одного раза при первом извлечении из очереди и не более одного раза при каждом возможном вытеснении, что даёт оценку общего числа операций с пулами, пропорциональную произведению числа абитуриентов на среднее число заявок одного абитуриента.
Операции вставки, удаления и чтения максимума в упорядоченном множестве выполняются за логарифмическое время от размера пула, а построение вектора оценок одного кандидата~--~за время, пропорциональное числу критериев в группе.

Итоговая временная сложность алгоритма выражается формулой~\ref{eq:time-complexity}:

\begin{equation}
\label{eq:time-complexity}
T = O(N \cdot M \cdot (\log K + R)),
\end{equation}

\begin{explanationx}
\item где $\mathrm{N}$~--~число абитуриентов;
\item $\mathrm{M}$~--~среднее число заявок одного абитуриента;
\item $\mathrm{K}$~--~максимальный размер пула категории, ограниченный её квотой;
\item $\mathrm{R}$~--~число критериев в группе категории.
\end{explanationx}

Затраты памяти определяются хранением очереди обработки, списков предпочтений абитуриентов и пулов категорий и выражаются формулой~\ref{eq:space-complexity}:

\begin{equation}
\label{eq:space-complexity}
S = O(N \cdot M).
\end{equation}

\subsubsection{Ранжирование кандидатов}

Компаратор кандидатов сравнивает их лексикографически по следующим параметрам в указанном порядке: вектор оценок по критериям категории по убыванию; значение \code{selectionpriority} заявки по возрастанию (более высокий приоритет предпочтения~--~меньшее число); идентификатор абитуриента по возрастанию (для обеспечения детерминированности и уникальности элементов упорядоченного множества).

Для критериев с вариантом \code{lower\_is\_better} значение инвертируется (умножается на минус единицу) перед сравнением, что позволяет использовать единый компаратор по принципу <<больше~--~лучше>>.

Блок-схема алгоритма представлена на рисунке~\ref{fig:selection-algo}.

\begin{figure}[htbp]
\centering
\includegraphics[width=\textwidth]{selection-algorithm}
\caption{Блок-схема алгоритма пересчёта результатов отбора}
\label{fig:selection-algo}
\end{figure}

\subsection{Архитектура подсистемы аудита и логического удаления}

\subsubsection{Журнал действий}

Подсистема аудита~--~единая точка записи действий пользователей, общая для всех репозиториев.
При выполнении любой операции записи репозиторий вызывает сервис аудита, который формирует запись в таблице \code{audit\_log}.
Состав полей таблицы приведён в таблице~\ref{tab:auditlog}.

\begin{longtblr}[
  caption = {Состав полей таблицы журнала действий},
  label = {tab:auditlog},
]{
  colspec = {|Q[l]|X|},
  width = \textwidth,
  hlines, vlines,
  rowhead = 1,
  row{1} = {halign=c},
}
Поле & Назначение \\
\code{id} & Первичный ключ записи \\
\code{user\_id} & Идентификатор пользователя, выполнившего действие \\
\code{username} & Имя пользователя на момент действия (на случай переименования) \\
\code{action} & Код действия \\
\code{entity\_type} & Тип изменённой сущности \\
\code{entity\_id} & Строковый идентификатор изменённой сущности \\
\code{changes} & Содержание изменения в формате JSONB \\
\code{created\_at} & Момент выполнения действия в часовом поясе UTC \\
\end{longtblr}

Поле \code{action} принимает одно из значений: \code{create} (создание), \code{update} (обновление), \code{delete} (запрос на удаление абитуриента или удаление иной сущности), \code{validate} (подтверждение валидации абитуриента), \code{invalidate} (снятие или автоматический сброс валидации), \code{delete\_confirmed} (подтверждение запроса на удаление), \code{delete\_rejected} (отклонение запроса на удаление), \code{recalculate\_all} (запуск пересчёта результатов отбора).

Поле \code{changes} содержит сериализованный объект с одним из трёх вариантов структуры: \code{\{after: \dots\}}~--~для создания (фиксируется состояние после изменения), \code{\{diff: \dots\}}~--~для обновления (фиксируется разница между состояниями до и после), \code{\{before: \dots\}}~--~для удаления (фиксируется состояние до изменения).
Такая структура позволяет восстановить любое состояние сущности по цепочке записей журнала.

\subsubsection{Сброс статуса валидации}

При любом изменении данных абитуриента сервис аудита автоматически сбрасывает поле \code{validated} в значение \code{false}, если оно ранее было установлено в \code{true}.
Одновременно в журнал записывается событие \code{invalidate} с признаком автоматического сброса, что позволяет отличать его от ручного снятия валидации пользователем.
Такое поведение реализует требование двойного контроля, сформулированное в подразделе~1.2.2: внесение и подтверждение данных выполняются двумя разными пользователями, а любое последующее изменение данных снова требует подтверждения.

Полный сценарий двойного контроля представлен на диаграмме деятельности, приведённой на рисунке~\ref{fig:activity-validation}.
На диаграмме выделены три дорожки: действия оператора приёмной комиссии, автоматические действия системы и действия аудитора.
Поток начинается с первичного ввода данных оператором, проходит через автоматическое сохранение с \code{validated = false} и пересчёт результатов отбора, затем переходит к проверке аудитором с принятием решения о подтверждении валидации.
Замыкающая ветвь показывает поведение системы при последующем изменении данных: автоматический сброс валидации и возврат записи на повторную проверку.

\begin{figure}[htbp]
\centering
\includegraphics[width=\textwidth]{activity-validation}
\caption{Диаграмма деятельности процесса двойного контроля данных абитуриента}
\label{fig:activity-validation}
\end{figure}

\subsubsection{Двухступенчатое удаление абитуриентов}

Двухступенчатое удаление абитуриентов основано на отдельной таблице \code{applicant\_deletion\_requests} с полем \code{status}, принимающим одно из трёх значений: \code{pending} (запрос подан, ожидает решения), \code{confirmed} (запрос подтверждён), \code{rejected} (запрос отклонён).
Возможные переходы между состояниями приведены в таблице~\ref{tab:deletionstates}.

\begin{longtblr}[
  caption = {Переходы состояний запроса на удаление абитуриента},
  label = {tab:deletionstates},
]{
  colspec = {|Q[l]|Q[l]|X|},
  width = \textwidth,
  hlines, vlines,
  rowhead = 1,
  row{1} = {halign=c},
}
Из состояния & В состояние & Условие перехода \\
(отсутствует) & \code{pending} & Пользователь с правом \code{WriteApplicants} нажал <<Удалить>> \\
\code{pending} & \code{confirmed} & Пользователь с правом \code{WriteAudit} подтвердил запрос \\
\code{pending} & \code{rejected} & Пользователь с правом \code{WriteAudit} отклонил запрос \\
\code{rejected} & (отсутствует) & Возможно создание нового запроса на того же абитуриента \\
\end{longtblr}

Абитуриент со статусом запроса \code{pending} или \code{confirmed} исключается из всех выборок репозитория абитуриентов и из снимка базы данных при пересчёте результатов отбора, что обеспечивает мгновенный эффект первичного запроса на удаление без потери данных.

При переводе запроса из состояния \code{pending} в \code{rejected} дополнительно выполняется сброс поля \code{validated} абитуриента в значение \code{false} и запуск повторного пересчёта результатов отбора с включением возвращённого в активные данные абитуриента.
Перевод запроса из состояния \code{pending} в \code{confirmed} не требует пересчёта, так как абитуриент уже был исключён из снимка при подаче запроса.
